\documentclass[12pt,a4paper]{article}

% =========================================================
% COMPILER OVERLEAF: XeLaTeX
% =========================================================
\usepackage{fontspec}
\usepackage[indonesian]{babel}

\IfFontExistsTF{Times New Roman}
    {\setmainfont{Times New Roman}}
    {\IfFontExistsTF{Times New Roman}
        {\setmainfont{Times New Roman}}
        {\setmainfont{Times New Roman}}}

% =========================================================
% PACKAGE
% =========================================================
\usepackage{graphicx}
\usepackage{float}
\usepackage{array}
\usepackage{tabularx}
\usepackage{longtable}
\usepackage{titlesec}
\usepackage{tocloft}
\usepackage{enumitem}
\usepackage{needspace}
\usepackage{ragged2e}
\usepackage{indentfirst}
\usepackage[none]{hyphenat}
\usepackage[hidelinks]{hyperref}

% =========================================================
% MARGIN SESUAI PEDOMAN
% =========================================================
\usepackage[
    a4paper,
    top=3cm,
    bottom=3cm,
    left=4cm,
    right=3cm
]{geometry}

% =========================================================
% PENGATURAN UMUM
% =========================================================
\setlength{\parindent}{1.25cm}
\setlength{\parskip}{0pt}

\setlength{\emergencystretch}{1em}

\hyphenpenalty=10000
\exhyphenpenalty=10000

\pretolerance=100
\tolerance=1000

\setlist[enumerate]{
    parsep=0pt,
    partopsep=0pt,
    listparindent=1.25cm
}

\newcommand{\headingfont}{\fontsize{14}{17}\selectfont}
\newcommand{\normalfontsize}{\fontsize{12}{18}\selectfont}

\newcommand{\placeholder}[1]{\textit{[#1]}}

\newcommand{\reportfigure}[4][0.88\textwidth]{%
    \begin{figure}[H]
        \centering
        \IfFileExists{#2}{%
            \includegraphics[
                width=#1,
                height=8cm,
                keepaspectratio
            ]{#2}%
        }{%
            \fbox{%
                \parbox[c][5.5cm][c]{0.82\textwidth}{%
                    \centering
                    File \texttt{#2} belum tersedia.\\
                    Masukkan tangkapan layar dengan nama file tersebut
                    ke folder yang sama dengan file LaTeX.
                }%
            }%
        }
        \caption{#3}
        \label{#4}
    \end{figure}
}

\setcounter{secnumdepth}{3}
\setcounter{tocdepth}{3}

\addto\captionsindonesian{%
    \renewcommand{\contentsname}{DAFTAR ISI}%
    \renewcommand{\listfigurename}{DAFTAR GAMBAR}%
    \renewcommand{\listtablename}{DAFTAR TABEL}%
    \renewcommand{\figurename}{Gambar}%
    \renewcommand{\tablename}{Tabel}%
}

% =========================================================
% FORMAT BAB DAN SUBBAB
% =========================================================
\renewcommand{\thesection}{\Roman{section}}
\renewcommand{\thesubsection}{\arabic{section}.\arabic{subsection}}
\renewcommand{\thesubsubsection}{\thesubsection.\arabic{subsubsection}}

\titleformat{\section}[display]
    {\centering\headingfont\bfseries}
    {BAB \thesection}
    {0pt}
    {\MakeUppercase}

\titlespacing*{\section}{0pt}{0pt}{1cm}

\titleformat{\subsection}
    {\headingfont\bfseries}
    {\thesubsection}
    {0.75em}
    {}

\titlespacing*{\subsection}{0pt}{0.8cm}{0.4cm}

\titleformat{\subsubsection}
    {\normalfontsize\bfseries}
    {\thesubsubsection}
    {0.75em}
    {}

\titlespacing*{\subsubsection}{0pt}{0.6cm}{0.25cm}

% =========================================================
% FORMAT DAFTAR ISI, GAMBAR, DAN TABEL
% =========================================================
\renewcommand{\cfttoctitlefont}{\hfill\headingfont\bfseries}
\renewcommand{\cftaftertoctitle}{\hfill\null}
\renewcommand{\cftloftitlefont}{\hfill\headingfont\bfseries}
\renewcommand{\cftafterloftitle}{\hfill\null}
\renewcommand{\cftlottitlefont}{\hfill\headingfont\bfseries}
\renewcommand{\cftafterlottitle}{\hfill\null}

\setlength{\cftbeforetoctitleskip}{0pt}
\setlength{\cftaftertoctitleskip}{1cm}
\setlength{\cftbeforeloftitleskip}{0pt}
\setlength{\cftafterloftitleskip}{1cm}
\setlength{\cftbeforelottitleskip}{0pt}
\setlength{\cftafterlottitleskip}{1cm}

\renewcommand{\cftsecpresnum}{BAB~}
\setlength{\cftsecnumwidth}{4.3em}
\setlength{\cftsecindent}{0pt}

\setlength{\cftsubsecindent}{1.2cm}
\setlength{\cftsubsecnumwidth}{1.3cm}

\setlength{\cftsubsubsecindent}{2.4cm}
\setlength{\cftsubsubsecnumwidth}{1.7cm}

\renewcommand{\cftsecleader}{\cftdotfill{\cftdotsep}}
\renewcommand{\cftsubsecleader}{\cftdotfill{\cftdotsep}}
\renewcommand{\cftsubsubsecleader}{\cftdotfill{\cftdotsep}}

\setlength{\cftbeforesecskip}{0.25cm}
\setlength{\cftbeforesubsecskip}{0pt}
\setlength{\cftbeforesubsubsecskip}{0pt}

\renewcommand{\cftsecfont}{\normalfont\normalfontsize}
\renewcommand{\cftsecpagefont}{\normalfont\normalfontsize}
\renewcommand{\cftsubsecfont}{\normalfont\normalfontsize}
\renewcommand{\cftsubsecpagefont}{\normalfont\normalfontsize}
\renewcommand{\cftsubsubsecfont}{\normalfont\normalfontsize}
\renewcommand{\cftsubsubsecpagefont}{\normalfont\normalfontsize}

\renewcommand{\cftfigfont}{\normalfont\normalfontsize}
\renewcommand{\cftfigpagefont}{\normalfont\normalfontsize}

\renewcommand{\cfttabfont}{\normalfont\normalfontsize}
\renewcommand{\cfttabpagefont}{\normalfont\normalfontsize}

% =========================================================
% AWAL DOKUMEN
% =========================================================
\begin{document}

\normalfontsize
\justifying

% =========================================================
% HALAMAN JUDUL / COVER
% =========================================================
\begin{titlepage}

    \thispagestyle{empty}
    \centering
    \setlength{\parindent}{0pt}

    {
        \headingfont
        \bfseries
        \addfontfeatures{LetterSpace=-3}

        \makebox[\textwidth][c]{%
            PENGEMBANGAN PASAR DIGITAL SEBAGAI MEDIA PROMOSI
        }\\[0.15cm]

        \makebox[\textwidth][c]{%
            PRODUK LOKAL BERBASIS WEB
        }\\[0.15cm]

        \makebox[\textwidth][c]{%
            DI PEKON SINAR MULYO
        }

        \par
    }

    \normalfontsize

    \vspace{2cm}

    {
        \bfseries
        PRAKTEK KERJA PENGABDIAN MASYARAKAT
        \par
    }

    \vspace{1.7cm}

    \IfFileExists{gambar/logo.png}
        {
            \includegraphics[width=5cm]{gambar/logo.png}
        }
        {
            \fbox{
                \parbox[c][5cm][c]{5cm}{
                    \centering
                    File logo.png\\
                    belum diunggah
                }
            }
        }

    \vspace{1.8cm}

    {
        \bfseries
        Disusun Oleh:
        \par
    }

    \vspace{0.4cm}

    {
        \bfseries

        \begin{tabular}{c@{\hspace{2cm}}c}
            Gerald Alpheratz Hussein & 2311010054
        \end{tabular}

        \par
    }

    \vfill

    {
        \bfseries
        FAKULTAS ILMU KOMPUTER\\
        PROGRAM STUDI TEKNIK INFORMATIKA\\
        INSTITUT INFORMATIKA DAN BISNIS DARMAJAYA\\
        BANDAR LAMPUNG\\
        2026
        \par
    }

\end{titlepage}

% =========================================================
% HALAMAN PENGESAHAN
% =========================================================
\clearpage

\pagenumbering{roman}
\setcounter{page}{1}

\thispagestyle{empty}
\setlength{\parindent}{0pt}

\phantomsection
\addcontentsline{toc}{section}{HALAMAN PENGESAHAN}

\begin{center}

    {
        \headingfont
        \bfseries
        HALAMAN PENGESAHAN
        \par
    }

    \normalfontsize

    \vspace{0.45cm}

    {
        \bfseries
        LAPORAN
        \par
    }

    \vspace{0.15cm}

    {
        \bfseries
        PRAKTEK KERJA PENGABDIAN MASYARAKAT (PKPM)
        \par
    }

    \vspace{0.75cm}

    \begin{minipage}{0.95\textwidth}

        \centering
        \bfseries

        PENGEMBANGAN PASAR DIGITAL SEBAGAI MEDIA PROMOSI\\
        PRODUK LOKAL BERBASIS WEB\\
        DI PEKON SINAR MULYO

    \end{minipage}

    \vspace{0.8cm}

    {
        \bfseries
        Oleh:
        \par
    }

    \vspace{0.4cm}

    {
        \bfseries

        \begin{tabular}{c@{\hspace{2.5cm}}c}
            Gerald Alpheratz Hussein & 2311010054
        \end{tabular}

        \par
    }

\end{center}

\vspace{1.1cm}

\begin{center}
    Telah memenuhi syarat untuk diterima.\\
    Menyetujui,
\end{center}

\vspace{0.7cm}

\noindent
\begin{tabularx}{\textwidth}{
    @{}
    >{\centering\arraybackslash}X
    >{\centering\arraybackslash}X
    @{}
}

    \textbf{Dosen Pembimbing}
    &
    \textbf{Pembimbing Lapangan}
    \\[2.4cm]

    \textbf{Lilla Rahmawati, S.Sos., M.M.}
    &
    \textbf{\placeholder{Nama Pembimbing Lapangan}}
    \\[0.15cm]

    NIK: \makebox[3.2cm][l]{1270213}
    &
    NIK: \makebox[3.2cm][l]{\dotfill}

\end{tabularx}

\vspace{1.4cm}

\begin{center}

    \textbf{Ketua Jurusan Teknik Informatika}

    \vspace{2.4cm}

    \textbf{Hary Sabita, S.T., M.T.I.}\\[0.15cm]

    NIK: \makebox[3.2cm][l]{\dotfill}

\end{center}

% =========================================================
% BAGIAN AWAL LAPORAN
% =========================================================
\clearpage

\pagestyle{plain}
\setlength{\parindent}{1.25cm}

% =========================================================
% DAFTAR ISI
% =========================================================
\phantomsection
\addcontentsline{toc}{section}{DAFTAR ISI}
\tableofcontents

% =========================================================
% DAFTAR GAMBAR
% =========================================================
\clearpage

\phantomsection
\addcontentsline{toc}{section}{DAFTAR GAMBAR}
\listoffigures

% =========================================================
% DAFTAR TABEL
% =========================================================
\clearpage

\phantomsection
\addcontentsline{toc}{section}{DAFTAR TABEL}
\listoftables

% =========================================================
% KATA PENGANTAR
% =========================================================
\clearpage

\phantomsection
\addcontentsline{toc}{section}{KATA PENGANTAR}

\begin{center}
    {
        \headingfont
        \bfseries
        KATA PENGANTAR
    }
\end{center}

\vspace{1cm}

Puji syukur penulis panjatkan ke hadirat Tuhan Yang Maha Esa atas
segala rahmat, karunia, dan petunjuk-Nya sehingga penulis dapat
menyelesaikan kegiatan serta menyusun Laporan Praktek Kerja
Pengabdian Masyarakat (PKPM) yang berjudul
\textbf{``Pengembangan Pasar Digital sebagai Media Promosi Produk
Lokal Berbasis Web di Pekon Sinar Mulyo''} dengan baik.

Laporan ini disusun sebagai bentuk pertanggungjawaban atas
pelaksanaan kegiatan PKPM sekaligus sebagai salah satu persyaratan
akademik di Institut Informatika dan Bisnis Darmajaya. Kegiatan
PKPM menjadi sarana bagi penulis untuk menerapkan ilmu pengetahuan
dan keterampilan yang telah diperoleh selama perkuliahan secara
langsung di tengah masyarakat.

Program pengembangan pasar digital Pekon Sinar Mulyo dilaksanakan
sebagai upaya mendukung pelaku Usaha Mikro, Kecil, dan Menengah
(UMKM) dalam memasarkan produk lokal secara lebih luas. Melalui
portal tersebut, informasi mengenai produk, harga, kategori, dan
lokasi usaha diharapkan dapat disampaikan secara lebih mudah, cepat,
dan terpercaya kepada calon pembeli.

Dalam pelaksanaan kegiatan dan penyusunan laporan ini, penulis
memperoleh bantuan, bimbingan, dukungan, dan kerja sama dari
berbagai pihak. Oleh karena itu, penulis menyampaikan terima kasih
kepada:

\begin{enumerate}[
    leftmargin=1.25cm,
    label=\arabic*.,
    itemsep=0pt,
    topsep=0.2cm
]

    \item Institut Informatika dan Bisnis Darmajaya yang telah
    memberikan kesempatan kepada penulis untuk melaksanakan
    kegiatan PKPM.

    \item \placeholder{Nama Ketua Jurusan} selaku Ketua Jurusan
    Teknik Informatika Institut Informatika dan Bisnis Darmajaya
    yang telah memberikan arahan dan dukungan dalam pelaksanaan
    kegiatan PKPM.

    \item \placeholder{Nama Dosen Pembimbing Lapangan} selaku Dosen
    Pembimbing Lapangan (DPL) yang telah memberikan bimbingan,
    arahan, dan masukan selama pelaksanaan kegiatan hingga
    penyusunan laporan ini.

    \item Kepala Pekon Sinar Mulyo beserta seluruh perangkat pekon
    yang telah memberikan izin, bantuan, informasi, dan dukungan
    selama pelaksanaan kegiatan PKPM.

    \item Para pelaku UMKM dan masyarakat Pekon Sinar Mulyo yang
    telah menerima, membantu, dan berpartisipasi dalam pelaksanaan
    kegiatan.

    \item Orang tua dan keluarga yang senantiasa memberikan doa,
    dukungan, serta motivasi kepada penulis.

    \item Rekan-rekan peserta PKPM yang telah bekerja sama dan
    saling membantu selama kegiatan berlangsung.

\end{enumerate}

Penulis menyadari bahwa laporan ini masih memiliki kekurangan,
baik dari segi penyusunan maupun isi. Oleh karena itu, penulis
mengharapkan kritik dan saran yang membangun sebagai bahan
perbaikan pada masa yang akan datang.

Penulis berharap laporan ini dapat memberikan manfaat bagi
Institut Informatika dan Bisnis Darmajaya, Pemerintah Pekon
Sinar Mulyo, para pelaku UMKM, serta pihak lain yang membutuhkan.
Semoga pasar digital yang dikembangkan dapat digunakan dan
dikembangkan secara berkelanjutan sebagai media promosi produk
lokal Pekon Sinar Mulyo.

\vspace{1.5cm}

\begin{flushright}

    Bandar Lampung,
    \makebox[3cm]{\dotfill}
    2026
    \\[1.8cm]

    \textbf{Gerald Alpheratz Hussein}

\end{flushright}

% =========================================================
% HALAMAN INTI
% =========================================================
\clearpage

\pagenumbering{arabic}
\setcounter{page}{1}
\setlength{\parindent}{1.25cm}

% =========================================================
% BAB I PENDAHULUAN
% =========================================================
\section{PENDAHULUAN}

\subsection{Latar Belakang}

Perekonomian desa sangat ditopang oleh aktivitas para pelaku Usaha
Mikro, Kecil, dan Menengah (UMKM). Di Pekon Sinar Mulyo terdapat
berbagai pelaku UMKM yang memproduksi produk lokal, antara lain
olahan makanan khas, hasil pertanian, kerajinan tangan, dan berbagai
jasa layanan. Produk-produk tersebut sebenarnya memiliki potensi dan
kualitas yang tidak kalah dengan produk dari daerah lain.

Namun, kondisi pemasaran produk lokal di Pekon Sinar Mulyo masih
dilakukan secara tradisional. Sebagian besar penjualan mengandalkan
pembeli yang datang langsung ke tempat usaha, atau promosi dari
mulut ke mulut yang jangkauannya sangat terbatas. Informasi produk
seperti harga, kontak penjual, dan lokasi tidak terdokumentasi
dengan baik. Akibatnya, produk lokal masih kesulitan bersaing dan
jangkauan pasarnya sempit.

Belum tersedianya media promosi yang terpusat juga menyebabkan
informasi mengenai produk UMKM Pekon Sinar Mulyo sulit ditemukan
oleh calon pembeli di luar wilayah pekon. Pembeli masih harus
memperoleh informasi produk secara langsung dari pelaku usaha atau
melalui komunikasi dari mulut ke mulut. Oleh karena itu, diperlukan
media resmi yang dapat menyajikan informasi produk secara
terstruktur, mudah dicari, dan dapat diakses kapan saja.

Salah satu solusi yang dapat diterapkan adalah mengembangkan portal
pasar digital berbasis web. Portal tersebut dapat digunakan untuk
menampilkan produk lokal, kategori produk, pencarian produk, harga,
lokasi usaha, hingga mekanisme verifikasi produk agar informasi yang
tampil kepada publik lebih terpercaya. Keberadaan portal diharapkan
dapat membantu pelaku UMKM memasarkan produk secara digital sekaligus
mempermudah pembeli menemukan produk yang dibutuhkan.

Portal yang dikembangkan tidak hanya menyediakan halaman produk
untuk pengunjung publik, tetapi juga dilengkapi sistem pengelolaan
dengan tiga peran pengguna, yaitu super admin, admin, dan pengguna
UMKM. Sistem tersebut memungkinkan pelaku UMKM memasukkan produk
secara mandiri, sedangkan admin bertugas memverifikasi produk
sebelum tampil kepada publik. Dengan demikian, pengelolaan produk
dapat dilakukan secara berkelanjutan dengan tetap menjaga kualitas
informasi yang disajikan.

\Needspace{7\baselineskip}

Melalui kegiatan Praktek Kerja Pengabdian Masyarakat (PKPM), penulis
menerapkan pengetahuan dan keterampilan di bidang teknologi
informasi dengan mengembangkan portal pasar digital bagi pelaku UMKM
Pekon Sinar Mulyo. Program ini diharapkan dapat mendukung
digitalisasi pemasaran, memperluas jangkauan produk lokal, serta
meningkatkan perekonomian masyarakat pekon.

\Needspace{8\baselineskip}

\subsubsection{Profil Pekon}

Pekon Sinar Mulyo merupakan salah satu pekon yang berada di
Kecamatan Pulau Panggung, Kabupaten Tanggamus, Provinsi Lampung.
Keberadaan Pekon Sinar Mulyo sebagai wilayah administrasi di
Kecamatan Pulau Panggung juga tercatat pada portal Kampung Keluarga
Berkualitas milik Kementerian Kependudukan dan Pembangunan Keluarga.
Secara administratif, Pekon Sinar Mulyo terbagi menjadi 2 dusun atau
Rukun Warga dan 6 Rukun Tetangga.

\reportfigure{gambar/pekon.png}
    {Dokumentasi Pekon Sinar Mulyo}
    {fig:profil-pekon-sinar-mulyo}

Berdasarkan data tahun 2024 yang diterbitkan oleh Badan Pusat
Statistik Kabupaten Tanggamus, Pekon Sinar Mulyo memiliki luas
wilayah 0,44 km\textsuperscript{2}. Jumlah penduduknya tercatat
sebanyak 1.048 jiwa yang terdiri atas 530 penduduk laki-laki dan
518 penduduk perempuan. Kepadatan penduduknya mencapai 2.381,82
jiwa per km\textsuperscript{2}. Pekon ini berjarak sekitar 5 km dari
ibu kota Kecamatan Pulau Panggung dan sekitar 23 km dari ibu kota
Kabupaten Tanggamus.

Kepadatan penduduk yang relatif tinggi dibandingkan dengan luas
wilayahnya menunjukkan besarnya jumlah rumah tangga yang berpotensi
menjalankan usaha mikro, kecil, dan menengah. Kondisi tersebut
memperkuat perlunya media promosi yang dapat menjangkau pembeli
di luar wilayah pekon, sehingga produk UMKM tidak hanya dipasarkan
secara terbatas di sekitar tempat tinggal pelaku usaha.

\noindent\textit{Sumber: Badan Pusat Statistik Kabupaten Tanggamus,
Kecamatan Pulau Panggung Dalam Angka 2025; serta Portal Kampung
Keluarga Berkualitas Kementerian Kependudukan dan Pembangunan
Keluarga.}

\Needspace{7\baselineskip}

\subsubsection{Potensi Ekonomi dan UMKM}

Potensi ekonomi Pekon Sinar Mulyo dapat dikembangkan lebih jauh
melalui penguatan pemasaran produk UMKM yang telah ada di wilayah
tersebut. Berdasarkan kondisi kelembagaan ekonomi dan jenis usaha
yang dijalankan masyarakat, potensi yang dapat dikembangkan melalui
pasar digital adalah sebagai berikut:

\begin{enumerate}[
    leftmargin=1.25cm,
    label=\arabic*.,
    itemsep=0pt,
    topsep=0.2cm
]

    \item \textbf{Produk olahan makanan}

    Pelaku UMKM memproduksi berbagai olahan makanan khas yang dapat
    dipasarkan lebih luas melalui kategori Makanan pada portal pasar
    digital.

    \item \textbf{Hasil pertanian}

    Kecamatan Pulau Panggung merupakan wilayah dengan aktivitas
    pertanian dan perkebunan, sehingga hasil pertanian pekon dapat
    dipromosikan melalui kategori Pertanian pada portal.

    \item \textbf{Kerajinan tangan}

    Produk kerajinan tangan hasil karya masyarakat dapat
    didokumentasikan dan dipasarkan melalui kategori Kerajinan.

    \item \textbf{Jasa layanan}

    Berbagai jasa layanan yang ditawarkan pelaku usaha lokal juga
    dapat dipromosikan melalui kategori Jasa pada portal, sehingga
    memperluas pilihan bagi calon pengguna layanan.

    \item \textbf{Konektivitas digital}

    Pekon Sinar Mulyo telah terjangkau jaringan 4G/LTE yang cukup
    kuat, sehingga mendukung penerapan pasar digital sebagai media
    promosi yang dapat diakses masyarakat maupun pembeli dari luar
    wilayah pekon.

\end{enumerate}

\noindent\textit{Sumber: Badan Pusat Statistik Kabupaten Tanggamus,
Kecamatan Pulau Panggung Dalam Angka 2025.}

\Needspace{8\baselineskip}

\subsection{Rumusan Masalah}

Berdasarkan latar belakang yang telah diuraikan, rumusan masalah
dalam pelaksanaan program ini adalah sebagai berikut:

\begin{enumerate}[
    leftmargin=1.25cm,
    label=\arabic*.,
    itemsep=0pt,
    topsep=0.2cm
]

    \item Bagaimana kondisi pemasaran produk lokal di Pekon Sinar
    Mulyo saat ini?

    \item Bagaimana merancang dan mengembangkan portal pasar digital
    berbasis web yang sesuai untuk mendukung pemasaran produk lokal?

    \item Bagaimana membangun sistem pengelolaan produk dengan
    mekanisme verifikasi agar informasi yang tampil kepada publik
    lebih terpercaya?

    \item Bagaimana implementasi dan dampak penerapan pasar digital
    terhadap pelaku UMKM dan masyarakat Pekon Sinar Mulyo?

\end{enumerate}

\Needspace{12\baselineskip}

\subsection{Tujuan}

Tujuan pelaksanaan program pengembangan pasar digital Pekon Sinar
Mulyo adalah sebagai berikut:

\begin{enumerate}[
    leftmargin=1.25cm,
    label=\arabic*.,
    itemsep=0pt,
    topsep=0.2cm
]

    \item Mengembangkan media promosi produk lokal berbasis web di
    Pekon Sinar Mulyo.

    \item Membantu pelaku UMKM memasarkan produk secara digital serta
    mempermudah pengelolaan data produk.

    \item Menyediakan sistem verifikasi produk dengan tiga peran
    pengguna, yaitu super admin, admin, dan pengguna UMKM, agar
    informasi produk yang tampil lebih terpercaya.

    \item Menerapkan pengetahuan dan keterampilan mahasiswa di bidang
    teknologi informasi melalui pelaksanaan kegiatan PKPM.

\end{enumerate}

\Needspace{12\baselineskip}

\subsection{Manfaat}

Manfaat yang diharapkan dari pelaksanaan program pengembangan pasar
digital Pekon Sinar Mulyo adalah sebagai berikut:

\begin{enumerate}[
    leftmargin=1.25cm,
    label=\alph*.,
    itemsep=0pt,
    topsep=0.2cm
]

    \item \textbf{Bagi Mitra (Pelaku UMKM dan Pemerintah Pekon)}

    Portal mempermudah promosi dan penjualan produk lokal sehingga
    jangkauan pasar lebih luas, sekaligus membantu perangkat pekon
    memfasilitasi pemasaran produk masyarakat.

    \item \textbf{Bagi Masyarakat}

    Masyarakat dan calon pembeli memperoleh kemudahan dalam mencari
    dan menemukan produk lokal Pekon Sinar Mulyo secara daring.

    \item \textbf{Bagi Mahasiswa}

    Program ini menjadi sarana untuk menerapkan ilmu yang dipelajari
    pada kondisi nyata di masyarakat, khususnya dalam perancangan dan
    pengembangan aplikasi web.

    \item \textbf{Bagi Institut Informatika dan Bisnis Darmajaya}

    Pelaksanaan program menjadi bentuk implementasi Tri Dharma
    Perguruan Tinggi melalui pengabdian kepada masyarakat.

\end{enumerate}

\Needspace{12\baselineskip}

\subsection{Teknologi yang Digunakan}

Pengembangan portal pasar digital Pekon Sinar Mulyo menggunakan
beberapa teknologi yang dipilih berdasarkan kebutuhan antarmuka
publik, pengelolaan produk, autentikasi peran pengguna, dan
penyimpanan data. Teknologi yang digunakan disajikan pada
Tabel~\ref{tab:teknologi-yang-digunakan}.

{
    \renewcommand{\arraystretch}{1.25}

    \begin{longtable}{
        |>{\raggedright\arraybackslash}p{2.7cm}
        |>{\raggedright\arraybackslash}p{3.8cm}
        |>{\raggedright\setlength{\parindent}{0pt}\arraybackslash}p{5.9cm}|
    }

        \caption{Teknologi yang digunakan dalam pengembangan portal}
        \label{tab:teknologi-yang-digunakan}\\

        \hline
        \textbf{Kategori}
        &
        \textbf{Teknologi}
        &
        \textbf{Kegunaan}
        \\
        \hline

        \endfirsthead

        \multicolumn{3}{c}{
            \tablename\ \thetable\ (lanjutan)
        }
        \\

        \hline
        \textbf{Kategori}
        &
        \textbf{Teknologi}
        &
        \textbf{Kegunaan}
        \\
        \hline

        \endhead

        \hline
        \multicolumn{3}{r}{
            Dilanjutkan pada halaman berikutnya
        }
        \\
        \endfoot

        \hline
        \endlastfoot

        Aplikasi web
        &
        React dan Vite
        &
        Digunakan untuk membangun halaman publik, halaman
        administrasi, komponen antarmuka, dan logika aplikasi
        portal pasar digital.
        \\
        \hline

        Basis data dan autentikasi
        &
        Supabase Auth, PostgreSQL, dan Supabase Storage
        &
        Digunakan untuk autentikasi tiga peran pengguna (super
        admin, admin, dan UMKM), penyimpanan data produk, serta
        penyimpanan gambar produk.
        \\
        \hline

        Pengelolaan hak akses
        &
        \textit{Row Level Security} (RLS)
        &
        Digunakan untuk membatasi akses data sesuai peran pengguna,
        sehingga pengguna UMKM hanya dapat mengelola produk miliknya
        sendiri sebelum diverifikasi admin.
        \\
        \hline

        Target publikasi
        &
        \placeholder{Nama layanan hosting}
        &
        Digunakan sebagai target publikasi aplikasi setelah
        konfigurasi, basis data, dan pemeriksaan akhir dinyatakan
        siap.
        \\

    \end{longtable}
}

Kombinasi teknologi tersebut mendukung pengembangan portal yang
responsif, terstruktur, dan mudah dikelola. Pemisahan akses antara
super admin, admin, dan pengguna UMKM diterapkan agar hanya produk
yang telah terverifikasi yang tampil kepada publik.

\Needspace{12\baselineskip}

\subsection{Mitra yang Terlibat}

Mitra yang terlibat dalam pelaksanaan program ini adalah pemerintah
Pekon Sinar Mulyo dan para pelaku UMKM setempat. Perangkat pekon
berperan sebagai fasilitator dan pemberi izin kegiatan, sedangkan
para pelaku UMKM berperan sebagai penyedia produk yang dipromosikan
serta calon pengguna sistem. Kerja sama ini menjadi dasar penting
dalam keberhasilan program pengembangan pasar digital.

% =========================================================
% BAB II PELAKSANAAN PROGRAM
% =========================================================
\clearpage

\section{PELAKSANAAN PROGRAM}

\subsection{Program-Program yang Dilaksanakan}

Program yang dilaksanakan dalam kegiatan Praktek Kerja Pengabdian
Masyarakat (PKPM) di Pekon Sinar Mulyo merupakan penerapan ilmu dan
keterampilan mahasiswa di bidang teknologi informasi. Program
disusun berdasarkan kebutuhan akan media promosi produk lokal yang
dapat mempertemukan pelaku UMKM dengan pembeli secara lebih luas.

Program utama yang dilaksanakan adalah pengembangan portal pasar
digital berbasis web beserta sistem verifikasi produk. Portal tidak
hanya dirancang sebagai halaman produk untuk pengunjung publik,
tetapi juga dilengkapi dashboard pengelolaan agar produk dapat
diperbarui secara mandiri oleh pelaku UMKM sebelum diverifikasi
admin. Uraian program yang dilaksanakan disajikan pada
Tabel~\ref{tab:program-yang-dilaksanakan}.

\Needspace{7\baselineskip}

\subsubsection{Program Kerja Utama}

Program kerja utama dalam kegiatan ini adalah pengembangan portal
pasar digital sebagai media promosi produk lokal berbasis web di
Pekon Sinar Mulyo.

{
    \renewcommand{\arraystretch}{1.2}

    \begin{longtable}{
        |>{\centering\setlength{\parindent}{0pt}\arraybackslash}p{0.7cm}
        |>{\raggedright\setlength{\parindent}{0pt}\arraybackslash}p{2.7cm}
        |>{\raggedright\setlength{\parindent}{0pt}\arraybackslash}p{5.7cm}
        |>{\raggedright\setlength{\parindent}{0pt}\arraybackslash}p{2.4cm}|
    }

        \caption{Program pengembangan pasar digital Pekon Sinar Mulyo}
        \label{tab:program-yang-dilaksanakan}\\

        \hline
        \textbf{No.}
        &
        \textbf{Kegiatan}
        &
        \textbf{Uraian}
        &
        \textbf{Sasaran}
        \\
        \hline

        \endfirsthead

        \multicolumn{4}{c}{
            \tablename\ \thetable\ (lanjutan)
        }
        \\

        \hline
        \textbf{No.}
        &
        \textbf{Kegiatan}
        &
        \textbf{Uraian}
        &
        \textbf{Sasaran}
        \\
        \hline

        \endhead

        \hline
        \multicolumn{4}{r}{Dilanjutkan pada halaman berikutnya}
        \\
        \endfoot

        \hline
        \endlastfoot

        1
        &
        Survei dan analisis kebutuhan
        &
        Melakukan survei lokasi serta analisis kebutuhan mitra untuk
        mengetahui produk lokal yang ada beserta kendala pemasarannya.
        &
        Pemerintah pekon dan pelaku UMKM.
        \\
        \hline

        2
        &
        Perancangan sistem
        &
        Menyusun struktur data, alur pengguna, pembagian tiga peran
        (super admin, admin, UMKM), serta rancangan tampilan
        antarmuka portal.
        &
        Super admin, admin, dan pengguna UMKM.
        \\
        \hline

        3
        &
        Pembangunan sistem
        &
        Membangun portal pasar digital menggunakan React dan Vite
        dengan \textit{backend} Supabase, meliputi halaman produk,
        kategori, pencarian, form tambah produk, dan dashboard
        verifikasi.
        &
        Pelaku UMKM dan pengunjung portal.
        \\
        \hline

        4
        &
        Pelatihan penggunaan sistem
        &
        Memberikan pelatihan kepada mitra agar dapat memasukkan dan
        mengelola produk secara mandiri melalui portal.
        &
        Pelaku UMKM.
        \\
        \hline

        5
        &
        Evaluasi dan penyempurnaan
        &
        Melakukan evaluasi kegiatan serta penyempurnaan sistem
        berdasarkan masukan mitra selama masa uji coba.
        &
        Pemerintah pekon dan pelaku UMKM.
        \\

    \end{longtable}
}

\Needspace{8\baselineskip}

\subsubsection{Program Kerja Tambahan}

\placeholder{
    Isi kegiatan tambahan jika ada, misalnya sosialisasi penggunaan
    portal kepada pelaku UMKM, gotong royong, atau kegiatan
    pendukung lain di luar pengembangan sistem inti.
}

\Needspace{8\baselineskip}

\subsubsection{Tahapan Pelaksanaan Program}

Pelaksanaan program pengembangan portal dilakukan melalui beberapa
tahap yang saling berkaitan. Tahap pertama berupa survei dan
identifikasi kebutuhan untuk menentukan produk serta kategori yang
perlu disediakan. Tahap berikutnya meliputi perancangan struktur
data, alur verifikasi, dan antarmuka portal. Setelah rancangan
disusun, pengembangan dilanjutkan pada halaman publik dan sistem
pengelolaan produk.

Tahap akhir meliputi pelatihan penggunaan sistem kepada pelaku UMKM,
pengujian fungsi verifikasi, serta penyempurnaan berdasarkan masukan
mitra. Tahapan tersebut dilakukan agar portal tidak hanya dapat
menampilkan produk, tetapi juga dapat dikelola secara mandiri dan
berkelanjutan oleh pelaku UMKM.

\subsection{Waktu Kegiatan}

Pelaksanaan kegiatan berlangsung selama periode pelaksanaan PKPM.
Jadwal kegiatan dirinci dalam Tabel~\ref{tab:waktu-pelaksanaan-bab-dua}.

{
    \renewcommand{\arraystretch}{1.2}

    \begin{longtable}{
        |>{\centering\setlength{\parindent}{0pt}\arraybackslash}p{0.8cm}
        |>{\raggedright\setlength{\parindent}{0pt}\arraybackslash}p{3cm}
        |>{\raggedright\setlength{\parindent}{0pt}\arraybackslash}p{7.4cm}|
    }

        \caption{Rangkaian waktu pelaksanaan kegiatan}
        \label{tab:waktu-pelaksanaan-bab-dua}\\

        \hline
        \textbf{No.}
        &
        \textbf{Tanggal}
        &
        \textbf{Kegiatan}
        \\
        \hline

        \endfirsthead

        \multicolumn{3}{c}{
            \tablename\ \thetable\ (lanjutan)
        }
        \\

        \hline
        \textbf{No.}
        &
        \textbf{Tanggal}
        &
        \textbf{Kegiatan}
        \\
        \hline

        \endhead

        \hline
        \multicolumn{3}{r}{Dilanjutkan pada halaman berikutnya}
        \\
        \endfoot

        \hline
        \endlastfoot

        1
        &
        \textit{23 juli 2026}
        &
        Survei lokasi dan analisis kebutuhan mitra.
        \\
        \hline

        2
        &
        \placeholder{isi tanggal}
        &
        Perancangan sistem dan antarmuka portal.
        \\
        \hline

        3
        &
        \placeholder{isi tanggal}
        &
        Pembangunan sistem portal pasar digital.
        \\
        \hline

        4
        &
        \placeholder{isi tanggal}
        &
        Pelatihan penggunaan sistem kepada mitra.
        \\
        \hline

        5
        &
        \placeholder{isi tanggal}
        &
        Evaluasi kegiatan dan penyempurnaan sistem.
        \\

    \end{longtable}
}

\subsection{Hasil Kegiatan dan Dokumentasi}

Berdasarkan program yang telah dilaksanakan, dihasilkan sebuah
portal pasar digital Pekon Sinar Mulyo yang dapat diakses melalui
browser. Portal menampilkan produk unggulan yang dikelompokkan ke
dalam kategori Makanan, Kerajinan, Pertanian, dan Jasa, serta
dilengkapi sistem verifikasi dengan tiga peran pengguna, yaitu super
admin, admin, dan pengguna UMKM.

\Needspace{8\baselineskip}

\subsubsection{Hasil Pengembangan Portal Publik}

Portal publik dirancang responsif agar dapat digunakan melalui
komputer, tablet, maupun telepon seluler. Pengunjung dapat mencari
produk, melihat detail harga dan satuan, memesan langsung kepada
penjual, serta membuka peta lokasi usaha.

\begin{enumerate}[
    leftmargin=1.25cm,
    label=\arabic*.,
    itemsep=0pt,
    topsep=0.2cm
]

    \item \textbf{Halaman Beranda}

    Halaman beranda menampilkan produk unggulan, kategori produk,
    dan fitur pencarian sebagai pintu masuk utama bagi pengunjung
    portal.

    \reportfigure{gambar/beranda.png}
        {Tampilan halaman beranda portal pasar digital}
        {fig:beranda-pasar}

    \item \textbf{Halaman Kategori Produk}

    Produk dikelompokkan ke dalam kategori Makanan, Kerajinan,
    Pertanian, dan Jasa agar pengunjung lebih mudah menemukan produk
    sesuai kebutuhan.

    \reportfigure{gambar/kategori.png}
        {Tampilan halaman kategori produk}
        {fig:kategori-produk}

    \item \textbf{Halaman Detail Produk}

    Setiap produk menampilkan informasi harga, satuan, deskripsi,
    kontak penjual, dan lokasi usaha pada peta.

    \reportfigure{gambar/detail.png}
        {Tampilan halaman detail produk}
        {fig:detail-produk}

    \item \textbf{Formulir Tambah Produk}

    Pengguna UMKM dapat memasukkan produk baru melalui formulir
    tambah produk yang kemudian menunggu proses verifikasi admin
    sebelum tampil kepada publik.

    \reportfigure{gambar/tambah.png}
        {Tampilan formulir tambah produk oleh pengguna UMKM}
        {fig:form-tambah-produk}

\end{enumerate}

\Needspace{8\baselineskip}

\subsubsection{Hasil Pengembangan Sistem Verifikasi dan Dashboard}

Sistem dibangun dengan tiga peran pengguna agar pengelolaan produk
berjalan terstruktur dan produk yang tampil kepada publik telah
melalui proses verifikasi.

\begin{enumerate}[
    leftmargin=1.25cm,
    label=\arabic*.,
    itemsep=0pt,
    topsep=0.2cm
]

    \item \textbf{Dashboard pengguna UMKM}

    Pengguna UMKM dapat menambahkan, menyunting, dan memantau status
    produk miliknya, baik yang masih menunggu verifikasi maupun yang
    sudah tampil di portal.

    \reportfigure{dashboard-umkm.png}
        {Tampilan dashboard pengguna UMKM}
        {fig:dashboard-umkm}

    \item \textbf{Dashboard admin}

    Admin bertugas memeriksa dan memverifikasi produk yang diajukan
    pengguna UMKM sebelum produk tersebut tampil kepada publik.

    \reportfigure{dashboard-admin-verifikasi.png}
        {Tampilan dashboard verifikasi produk oleh admin}
        {fig:dashboard-admin-verifikasi}

    \item \textbf{Dashboard super admin}

    Super admin memiliki kewenangan tertinggi, meliputi pengelolaan
    akun admin, pemantauan seluruh data produk, dan pengaturan umum
    portal.

\end{enumerate}

\Needspace{8\baselineskip}

\subsubsection{Penerapan Basis Data dan Keamanan}

Data portal disimpan menggunakan PostgreSQL melalui Supabase. Tabel
utama meliputi data akun pengguna, produk, kategori, dan status
verifikasi. Setiap tabel dilindungi menggunakan \textit{Row Level
Security} (RLS) agar pengguna UMKM hanya dapat mengelola produk
miliknya sendiri, sedangkan admin dan super admin memiliki
kewenangan sesuai peran masing-masing.

\subsection{Dampak Kegiatan}

Program pengembangan pasar digital Pekon Sinar Mulyo memberikan
hasil berupa fondasi media promosi produk lokal berbasis web.
Dampak program tidak hanya berkaitan dengan tersedianya portal,
tetapi juga dengan perubahan cara produk UMKM dipasarkan dan
dikelola.

\Needspace{7\baselineskip}

\subsubsection{Dampak bagi Pelaku UMKM}

Pelaku UMKM memperoleh media tertata untuk mengelola data produk,
sehingga informasi harga, kontak, dan lokasi tersusun dalam satu
tempat. Jangkauan pemasaran menjadi lebih luas karena produk dapat
dikenal masyarakat di luar pekon, dan proses persuasi kepada calon
pembeli menjadi lebih terpercaya berkat mekanisme verifikasi.

\Needspace{7\baselineskip}

\subsubsection{Dampak bagi Masyarakat dan Pembeli}

Masyarakat dan calon pembeli memperoleh kemudahan dalam mencari
produk lokal Pekon Sinar Mulyo tanpa harus datang langsung ke tempat
usaha. Tampilan responsif memudahkan penggunaan portal melalui
telepon seluler kapan saja dibutuhkan.

\Needspace{7\baselineskip}

\subsubsection{Dampak bagi Mahasiswa}

Pelaksanaan program memberikan pengalaman kepada mahasiswa dalam
menganalisis kebutuhan mitra, merancang sistem tiga peran pengguna,
mengembangkan antarmuka responsif, mengelola basis data dan
autentikasi, serta melakukan pelatihan penggunaan sistem kepada
masyarakat.

% =========================================================
% BAB III PENUTUP
% =========================================================
\clearpage

\section{PENUTUP}

\subsection{Kesimpulan}

Kegiatan PKPM ini berhasil mengembangkan portal pasar digital
berbasis web di Pekon Sinar Mulyo sebagai media promosi produk
lokal. Sistem tersebut mempertemukan pelaku UMKM dengan pembeli,
memudahkan pengelolaan data produk, serta dilengkapi dengan
mekanisme verifikasi agar informasi yang tampil lebih terpercaya.
Melalui kegiatan ini, masyarakat setempat mendapatkan wadah digital
untuk memperluas pemasaran produk lokal.

\subsection{Saran}

Beberapa saran untuk keberlanjutan program adalah sebagai berikut.

\begin{enumerate}[
    leftmargin=1.25cm,
    label=\arabic*.,
    itemsep=0pt,
    topsep=0.2cm
]

    \item Bagi mitra: agar rutin memperbarui data produk dan
    memanfaatkan sistem secara konsisten.

    \item Bagi mahasiswa berikutnya: dapat menambahkan fitur
    pembayaran dan laporan penjualan.

    \item Bagi kampus: agar terus mendukung kegiatan pengabdian yang
    berkelanjutan.

\end{enumerate}

\subsection{Rekomendasi}

Keberlanjutan program sebaiknya dilakukan pendampingan lanjutan agar
seluruh pelaku UMKM mahir menggunakan sistem. Dinas atau perangkat
desa dapat mengelola data produk dan menjadikan portal sebagai media
resmi untuk mendukung pertumbuhan ekonomi desa. Dengan demikian,
manfaat pasar digital dapat terus dirasakan dalam jangka panjang.

% =========================================================
% DAFTAR PUSTAKA
% =========================================================
\clearpage

\phantomsection
\addcontentsline{toc}{section}{DAFTAR PUSTAKA}

\begin{center}
    {
        \headingfont
        \bfseries
        DAFTAR PUSTAKA
    }
\end{center}

\vspace{1cm}

\begin{flushleft}
Badan Pusat Statistik Kabupaten Tanggamus. (2025). \textit{Kecamatan
Pulau Panggung Dalam Angka 2025}. BPS Kabupaten Tanggamus.

\placeholder{Tambahkan referensi lain yang digunakan, misalnya
jurnal, buku, atau dokumentasi teknis React/Supabase, dengan format
Nama Penulis. (Tahun). Judul Referensi. Nama Jurnal/Penerbit.}
\end{flushleft}

% =========================================================
% LAMPIRAN-LAMPIRAN
% =========================================================
\clearpage

\phantomsection
\addcontentsline{toc}{section}{LAMPIRAN-LAMPIRAN}

\begin{center}
    {
        \headingfont
        \bfseries
        LAMPIRAN-LAMPIRAN
    }
\end{center}

\vspace{1cm}

\textbf{
    Lampiran 1. Bukti Kegiatan Selama Pelaksanaan PKPM
}

\vspace{0.5cm}

\placeholder{
    Masukkan foto, dokumen, daftar hadir,
    atau bukti kegiatan selama pelaksanaan PKPM.
}

\vspace{1cm}

\textbf{
    Lampiran 2. Bukti Pendukung Program yang Dilaksanakan
}

\vspace{0.5cm}

\placeholder{
    Masukkan bukti pendukung program,
    misalnya tangkapan layar portal pasar digital,
    hasil serah terima, surat,
    atau dokumen pendukung lainnya.
}

\end{document}